\section{Constructing the quadratic ray-class family}
\label{sec:quadratic-ray-family}

Let \(d\geq100\) be a prime with \(d\equiv3\bmod4\), let \(M\geq3\),
and set
\[
 K=\mathbb Q(\sqrt{-d}),\quad
 {\cal O}=\mathbb Z[\omega],\quad
 \omega=(1+\sqrt{-d})/2,\quad D=4dM^2.
\]
The discriminant is \(-d\), and the area of \({\cal O}\) in
\(\mathbb C\) is \(\sqrt d/2\).
The norm formula
\(N((A+B\sqrt{-d})/2)=(A^2+dB^2)/4\) shows that the units are
\(1,-1\). Since \(M\geq3\), the only unit congruent to one modulo
\(M{\cal O}\) is one.

Let \(G\) be the ray class group for \(M{\cal O}\), with no infinite
part. Its elements are fractional ideals prime to \(M\), modulo
principal ideals generated by elements congruent to one at the primes
of \(M\). Write its law additively and its trivial class as zero.
Every class has an integral representative prime to \(M\): given a
fractional representative, choose a denominator congruent to one at
the modulus and divisible by all its other denominator primes, using
the Chinese remainder theorem. Multiplication by that denominator
leaves the ray class unchanged.

The ray-class exact sequence gives
\[
 |G|\leq h_K|({\cal O}/M{\cal O})^\times|\leq h_KM^2.
\]
For completeness, the first inequality follows by mapping an ideal
to its ordinary ideal class. Its kernel is generated by principal
ideals prime to \(M\), whose ray classes are determined by the
generator's residue modulo \(M\), up to a unit. Reduction of a
fractional generator uses its local denominator inverse at \(M\).
Every ordinary ideal class has an ideal prime to \(M\) by the same
approximation argument.

Minkowski's ideal-class bound in this quadratic imaginary field is
\((2/\pi)\sqrt d\). The number of ideals of norm \(n\) is at most
the divisor function \(\tau(n)\). This follows prime by prime: a
split rational prime gives \(k+1\) ideals of norm \(p^k\), an inert
one gives zero or one, and a ramified one gives one. Hence, putting
\(Y=(2/\pi)\sqrt d\),
\[
 h_K\leq\sum_{n\leq Y}\tau(n)
       =\sum_{ab\leq Y}1
       \leq Y(1+\log Y)\leq4d .
\]
The last inequality follows from \(Y\leq\sqrt d\),
\(1+\log\sqrt d\leq4\sqrt d\), and \(d\geq100\).
Thus \(G\) is finite and
\begin{equation}
 \label{eq:ray-group-budget}
                            H:=|G|\leq D.
\end{equation}

\subsection{One lattice for each ray class}

Let \(a_g(n)\) count integral ideals prime to \(M\), in class \(g\),
of norm \(n\). Fix an integral ideal \(J\) prime to \(M\) representing
\(-g\). The Chinese remainder theorem supplies
\(\alpha_0\in J\) with \(\alpha_0\equiv1\bmod M{\cal O}\).
The map
\begin{equation}
 \label{eq:ray-lattice-bijection}
 \alpha\longmapsto(\alpha)J^{-1}
\end{equation}
is a bijection from \(\alpha_0+MJ\) to the ideals counted in
\(a_g\). An element in this coset lies in \(J\), so its image is
integral. It is prime to \(M\), has ray class \(g\), and is nonzero.
Conversely \(I\) in class \(g\) makes \(IJ\) ray principal, with
a generator \(\alpha\equiv1\bmod M\). This generator lies in \(J\).
Also
\(J\cap M{\cal O}=MJ\) by coprimality, so it lies in the indicated
coset. Two such generators of the same ideal differ by a unit
congruent to one modulo \(M\), which is one. This proves injectivity.
The assertion about an integral generator congruent to one follows
also with the localization definition of the ray relation. An integral
element of \(K\) congruent to one at every prime of \(M\) has difference
from one in \(M{\cal O}\), as is seen from the valuations at those primes.

Use the real quadratic form
\[
                         q(\alpha)=\frac{|\alpha|^2}{N(J)}
\]
on this coset. For \(\alpha\in J\), its value is the norm of
\((\alpha)J^{-1}\), and is therefore a positive integer on
\(\alpha_0+MJ\).

Normalize the lattice \(MJ\subset\mathbb C\) by \(N(J)^{-1/2}\).
Its area is \(M^2\sqrt d/2\).
Every nonzero vector is \(M\beta/N(J)^{1/2}\) with \(\beta\in J\),
and its squared length is at least \(M^2\), since
\(N((\beta)J^{-1})\) is a positive integer.
Choose a shortest nonzero vector \(u\). It is primitive, since an
integer multiple of another lattice vector would be longer.
Extend it to a basis \(u,v\), and replace \(v\) by \(v-ku\), with
\(k\) the nearest integer to its projection coefficient.
Then, writing
\[
 \alpha=|u|^2,\qquad
 l=\frac{\langle u,v\rangle}{|u|^2},\qquad
 \gamma=|v|^2-l^2|u|^2,
\]
we have
\[
 \alpha\geq M^2,\quad |l|\leq1/2,\quad
 \gamma\geq3\alpha/4\geq1,\quad
                   \alpha\gamma=dM^4/4 .
\]
Indeed \(|v|\geq|u|\) by minimality of \(u\).
In these coordinates the coset has form
\[
 Q(m,n)=\alpha\{m+\theta_1+l(n+\theta_2)\}^2
                           +\gamma(n+\theta_2)^2 .
\]
It follows from \eqref{eq:ray-lattice-bijection} that
\begin{equation}
 \label{eq:ray-fiber-count}
        a_g(k)=\#\{(m,n)\in\mathbb Z^2:Q(m,n)=k\}.
\end{equation}
All values \(Q(m,n)\) are positive integers. The constants
\(\theta_1,\theta_2,l,\alpha,\gamma\) may depend on \(g\), but
their product \(\alpha\gamma\) is the same for every class.

\subsection{Uniform counting and theta estimates}

For \(x\geq1\), a row \(n+\theta_2=y\) contributes integers \(m\)
in an interval of length
\[
                 f(y)=2\sqrt{(x-\gamma y^2)_+/\alpha}.
\]
For rows with \(|y|\leq\sqrt{x/\gamma}\), the discrepancy between
their integer count and that length is at most one.
There are at most \(2\sqrt{x/\gamma}+1\) such rows.
The function \(f\) is nonnegative and unimodal, with total variation
\(2\sup f=4\sqrt{x/\alpha}\).
Comparing its value at the left endpoint of each interval of length
one with its integral on that interval gives
\[
       \left|\sum_{n\in\mathbb Z}f(n+\theta_2)
                              -\int_{\mathbb R}f(y)\,dy\right|
                           \leq4\sqrt{x/\alpha}.
\]
To see this inequality directly, the discrepancy on one interval is
at most the variation of \(f\) on that interval. Sum these variations
over the partition. The function has compact support, so the sum is
its total variation.
The ellipse area is
\(\int f=\pi x/\sqrt{\alpha\gamma}\). Therefore
\[
 \left|\#\{(m,n):Q(m,n)\leq x\}
                         -\frac{\pi x}{\sqrt{\alpha\gamma}}\right|
 \leq2\sqrt{x/\gamma}+1+4\sqrt{x/\alpha}
                                  \leq10\sqrt x.
\]
Together with \eqref{eq:ray-fiber-count}, this proves
\begin{equation}
 \label{eq:ray-count-done}
 \left|\sum_{n\leq x}a_g(n)-\kappa x\right|\leq10\sqrt x,
                        \qquad \kappa=\frac{2\pi}{\sqrt d M^2}.
\end{equation}

We also prove the theta estimate uniformly in the shifts and shear.
For \(a>0\), the periodized Gaussian has Fourier series
\[
 \sum_{m\in\mathbb Z}e^{-a(m+c)^2}
       =\sqrt{\frac{\pi}{a}}\sum_{k\in\mathbb Z}
                         e^{-\pi^2k^2/a}e^{2\pi ikc}.
\]
Its coefficients are obtained by integrating on one period and
translating to the Gaussian integral on \(\mathbb R\).
Both the periodized sum and its Fourier series converge absolutely.
The Gaussian Fourier integral follows by differentiation in frequency
and the Gaussian integral at zero.
Since \(e^{-z}\leq2z^{-2}\) for \(z>0\), and
\(\sum_{k\geq1}k^{-4}\leq4/3\), this implies
\[
 \left|\sum_me^{-a(m+c)^2}-\sqrt{\pi/a}\right|
   \leq\sqrt{\pi/a}\frac{4a^2}{\pi^4}\sum_{k\geq1}k^{-4}
   \leq\sqrt{\pi/a}\,a^2 .
\]
The last inequality uses \(\pi>3\).

Apply this first in \(m\), at \(a=\alpha t\), uniformly in
\(c=\theta_1+l(n+\theta_2)\), and then in \(n\), at
\(a=\gamma t\). Positivity of the outer Gaussian gives
\[
 \left|\sum_{m,n}e^{-tQ(m,n)}
                      -\frac{\pi}{t\sqrt{\alpha\gamma}}\right|
 \leq\frac{\pi}{t\sqrt{\alpha\gamma}}
        \{(\alpha t)^2+(\gamma t)^2+(\alpha t)^2(\gamma t)^2\}.
\]
The same estimates establish absolute convergence of the double
sum. For \(0<t\leq1\), \(\alpha,\gamma\geq1\) imply
\[
 \alpha^2+\gamma^2+\alpha^2\gamma^2t^2
                                     \leq3(\alpha\gamma)^2.
\]
Since \(\alpha\gamma=dM^4/4\), the error is at most
\[
        3\pi(\alpha\gamma)^2t
             =\frac{3\pi}{16}d^2M^8t\leq D^{10}t .
\]
Grouping the positive terms by their positive integer norm
\(Q(m,n)\) gives the desired identity
\begin{equation}
 \label{eq:ray-theta-done}
       \left|\sum_{n\geq1}a_g(n)e^{-tn}-\frac{\kappa}{t}\right|
                                  \leq D^{10}t .
\end{equation}
The constants also satisfy
\(D^{-1}\leq\kappa\leq4\), by \(D=4dM^2\), \(d\geq100\),
\(M\geq3\) and \(3<\pi<4\).

\subsection{The prime weights and the positive convolution}

For an integral ideal \(I\) prime to \(M\), define
\(\Lambda_K(I)=\log N(P)\) when \(I=P^j\) for a prime ideal \(P\)
and \(j\geq1\), and zero otherwise. Define
\[
       \lambda_g(n)=
        \sum_{\substack{N(I)=n\\ [I]=g\\(I,M)=1}}\Lambda_K(I).
\]
These weights are nonnegative. The only nonzero terms occur at
rational prime powers \(n=p^k\).
For a prime \(P\) of residue degree \(f_P\) over \(p\), a contribution
at \(p^k\) is \(f_P\log p\), and occurs only if \(f_P\mid k\).
Since \(\sum_{P\mid p}e_Pf_P=2\), we get
\begin{equation}
 \label{eq:ray-lambda-budget}
                         \sum_g\lambda_g(n)\leq2\Lambda(n).
\end{equation}
Omitting the primes dividing the modulus can only decrease the sum.

Unique factorization of ideals gives
\(\log N(I)=\sum_{J\mid I}\Lambda_K(J)\).
Partitioning the divisors by norm and ray class proves
\[
 a_g(n)\log n
       =\sum_{h\in G}\sum_{ab=n}\lambda_h(a)a_{g-h}(b).
\]
All sums are finite. Divisors and complementary divisors of an ideal
prime to \(M\) are also prime to \(M\).

For a real character \(\chi\), the ideal coefficient of the
convolution of the trivial and \(\chi\) series at an ideal
\(C=\prod_PP^{e_P}\) is
\[
             r(C)=\sum_{J\mid C}\chi([J])
                      =\prod_P\sum_{j=0}^{e_P}\chi([P])^j\geq0.
\]
Each factor is \(e_P+1\) when \(\chi([P])=1\), and is zero or
one when \(\chi([P])=-1\). Grouping by \(N(C)=n\) proves
\[
                      \sum_{ab=n}a_1(a)a_\chi(b)\geq0,
\]
which is \eqref{eq:ray-convolution-positive}.
The second positivity property in the formal family follows too.
For every ideal \(I\), all exponents in \(I^2\) are even, so
\(r(I^2)\geq1\). The map \(I\mapsto I^2\) is injective by unique
factorization, and \(N(I^2)=N(I)^2\).
It follows that
\[
                  a_1(n)\leq\sum_{ab=n^2}a_1(a)a_\chi(b).
\]

Finally, if \(p\) is prime and \(\lambda_0(p)\ne0\), some ideal of
norm \(p\) is ray principal. It has a generator
\(\alpha\equiv1\bmod M{\cal O}\).
Write
\(\alpha=1+M(a+b\omega)=(A+B\sqrt{-d})/2\).
Then
\[
  A=2+2Ma+Mb,\qquad B=Mb,\qquad A^2+dB^2=4p.
\]
Thus \(M\mid B\) and \(2M\mid A-2-B\), which is precisely
\(\operatorname{NormOne}(d,M,p)\).

We have constructed every counting, theta, convolution and support
property of the family used in
Sections~\ref{sec:ray-real-gap} and~\ref{sec:level-prime-selection},
with the same parameter \(D=4dM^2\). The remaining analytic input
to the norm-prime count is the characterwise explicit formula
with its unique possible exceptional zero.

\paragraph{Source map.}
This argument follows
\srcid{RayClassIdealCountEqNcardShiftedBinaryForm},
\srcid{QFormCount}, \srcid{ThetaQForm},
\srcid{RayCountThetaFibreCount}, \srcid{FiniteIdealsIn},
\srcid{ConvRealChar}, and \srcid{ExistsLFamRay}.
Minkowski's ideal-class bound and the ray-class exact sequence are
used in their explicit standard forms stated above.
\footnote{\href{https://github.com/Deicyde/riemann-hypothesis/blob/40e56bab026552e089d7afcfeff08d2defb9c005/src/Development/AbelianVarietyTheory/Proofs/GlobalInterpolation/ExistsLFamRay.lean}{Assembly of the ray family};
\href{https://github.com/Deicyde/riemann-hypothesis/blob/40e56bab026552e089d7afcfeff08d2defb9c005/src/Development/AbelianVarietyTheory/Proofs/GlobalInterpolation/RayClassIdealCountEqNcardShiftedBinaryForm.lean}{the ideal and lattice bijection};
\href{https://github.com/Deicyde/riemann-hypothesis/blob/40e56bab026552e089d7afcfeff08d2defb9c005/src/Development/AbelianVarietyTheory/Proofs/GlobalInterpolation/ConvRealChar.lean}{the two real-character positivity assertions}.}
